\label{chap:Introduction}
This thesis presents a Procedural Machine Learning Pipeline for Earthquake Pick 
Detection and Phase Association, integrating \ac{HPC} infrastructures to 
enhance both the accuracy and scalability of seismic data analysis. The work 
focuses on seismic waveforms recorded by the \ac{SMINO}, managed by the 
\ac{OGS}, in Northeastern Italy and surrounding regions including Veneto, 
Slovenia, and Alto Adige.

Earthquakes are caused by the sudden release of elastic strain energy 
accumulated within the Earth's lithosphere. This energy propagates as seismic 
waves, which travel through heterogeneous geological structures. The 
lithosphere constitutes a mechanically rigid yet dynamically active shell and 
its behavior produces the deformation patterns that manifest as seismicity and 
volcanism. Capturing the temporal and spatial variability of these processes 
requires robust and high throughput analytical tools.

Seismometers record ground motion as continuous digital time series, but the 
transformation of these raw signals into scientifically meaningful information 
depends on the accurate detection of seismic phases. Phase picking, the 
identification of P- and S-wave onset times, is foundational for earthquake 
location, source characterization, and downstream seismological analyses. 
Traditional automated picking algorithms such as the \ac{STA/LTA} 
(\cite{STALTA}) method and the \ac{AIC} picker have provided decades of 
utility. However, their performance degrades substantially under challenging 
conditions: low signal to noise ratios, emergent onsets, overlapping events, 
and environmental noise. Furthermore, operational seismic networks often depend 
on expert analysts to review and correct automated picks, introducing latency 
and potential subjective bias.

Continuous waveform archives from regional and global networks commonly reach 
terabyte scales annually, rendering manual processing infeasible. Addressing 
this "Big Data" problem requires distributed architectures employing parallel 
algorithms, GPU acceleration, and scalable storage. This thesis integrates 
\ac{ML} workflows with \ac{HPC} clusters to enable efficient processing at 
operational scale.

The inherent complexity and nonlinearity of seismic phenomena have motivated 
the adoption of machine learning approaches capable of discovering subtle 
waveform patterns beyond the reach of rule based algorithms. Recent advances in 
\ac{DL} have led to models that outperform classical methods in event 
detection, signal classification, and phase arrival estimation. These include 
convolutional and recurrent neural architectures trained on large labeled 
data sets, which demonstrate improved robustness to noise and generalization 
across stations and regions. Their ability to extract multi-scale temporal
features makes them particularly well suited to the intricate structure of 
seismic waveforms.

By coupling \ac{ML} methodologies with large scale computational frameworks, 
this thesis aims to refine pick detection and event association in Northeastern 
Italy. The regional focus is particularly relevant, as this area is 
characterized by complex tectonic structures and moderate to high seismicity. 
Improving monitoring capabilities here has direct implications for hazard 
assessment, early warning, and civil protection strategies. A noteworthy 
contribution to this field is SeisBench (\cite{SeisBench}), an open-source 
platform that offers unified access to state-of-the-art \ac{DL} models for 
seismic analysis. By providing standardized APIs, pretrained models, and 
curated data sets, SeisBench reduces the barriers to reproducibility and 
facilitates rapid experimentation. This thesis leverages SeisBench to access 
and evaluate pretrained PhaseNet and EQTransformer models trained on multiple 
benchmark datasets.

The focus of this work lies on seismic waveforms recorded by the \ac{SMINO}, a 
northeastern Italian network (\cite{SMINO}), established following the 
destructive 1976 Mw 6.4 Friuli earthquake. The network plays a critical role in 
real-time seismic surveillance and civil protection activities in the region. 
The thesis evaluates \ac{ML} based models for automatic pick detection, event 
association, hypocenter location, and magnitude estimation on this dataset, 
with an emphasis on deploying methods suitable for near real-time operation.

The evaluation is as follows: two deep-learning picker architectures (PhaseNet 
(\cite{PhaseNet}) and EQTransformer (\cite{EQTransformer})) are each trained on 
four benchmark datasets (INSTANCE (\cite{INSTANCE}), STEAD (\cite{STEAD}), 
SCEDC (\cite{SCEDC}), and the original training corpus), and their performance 
is assessed across three detection-probability thresholds (0.1, 0.2, 0.3). The 
best-performing picker is then paired with two association algorithms (GaMMA 
(\cite{GaMMA}) and PyOcto (\cite{PyOcto})) and a probabilistic locator, 
NonLinLoc (\cite{NonLinLoc}) to quantify end-to-end pipeline recall, timing 
precision, spatial accuracy, and magnitude fidelity against the OGS reference 
catalog. To compare the ML results with the OGS reference catalog, OGS events 
were relocated using the 1D NonLinLoc (\cite{NonLinLoc}) configuration adopted 
in this study. The catalog will be published soon (Magrin, personal 
communication). It contains 380 events within the red polygon shown in 
Figure~\ref{fig:OGSSeismicity}, with local magnitudes ranging from -0.2 to 4.6. 
A total of 7,129 phase picks were associated with these events, including 
4,126~P picks and 3,003~S picks.

The document is structured as follows. Chapter~2 introduces the data sets 
utilized in this study, detailing acquisition, metadata standards, and data 
quality control procedures. Chapter~3 describes the data engineering 
framework, encompassing multi-format catalog parsing, data aggregation, and 
visualization components. Chapter~4 presents the computational pipeline 
architecture, describing the preprocessing steps and the implementation of 
deep-learning models for phase picking, event association, hypocenter 
location, and magnitude estimation. Chapter~5 presents the pipeline evaluation 
framework, introducing a physically informed validation methodology based on 
bipartite graph matching and domain-specific confusion matrices to rigorously 
assess model performance against manual catalogs. Chapter~6 details the 
high-performance computing infrastructure and parallelization strategy deployed 
on the Leonardo supercomputer at CINECA, including job submission, resource 
management, and scalability considerations. Chapter~7 presents the results of 
the comprehensive benchmarking campaign on the Northeastern Italy dataset, 
analyzing performance through recall statistics, residual distributions, and 
detailed comparisons between models trained on different data sets. Finally, 
Chapter~8 synthesizes the key findings, discusses the implications for 
operational seismology, and outlines future research directions for scalable, 
automated monitoring systems.

\section{Seismic Waveforms}
An \textbf{earthquake} is the shaking of the Earth's surface resulting from a 
sudden release of energy in the lithosphere, which generates 
\textbf{seismic waves}. These events arise primarily from the rupture of 
geological faults under tectonic stress, ranging in magnitude from 
imperceptible events to catastrophic events capable of widespread destruction. 
Seismic waves propagate energy through the Earth and are broadly categorized 
into two main types: \textbf{body waves} and \textbf{surface waves}.

\begin{itemize}
  \item \textbf{Body waves} are seismic waves that travel through the Earth's 
  interior and are the fastest seismic waves. Body waves are further divided 
  into two types: \textbf{P waves} and \textbf{S waves}.
  \begin{itemize}
    \item \textbf{P waves} (Primary or Compressional waves): These are the 
    fastest seismic waves and the first to be recorded by seismometers. They 
    propagate through solids, liquids, and gases by compressing and expanding 
    the medium in the direction of propagation (longitudinal motion).
    \item \textbf{S waves} (Secondary or Shear waves): S waves are slower than 
    P waves and arrive second. They propagate only through solid materials, 
    moving the medium perpendicular to the direction of propagation (transverse 
    motion).
  \end{itemize}
  \item \textbf{Surface waves} propagate along the Earth's surface. While 
  slower than body waves, they typically have larger amplitudes and are 
  responsible for the majority of structural damage. They include \textbf{Love 
  waves} (horizontal shear) and \textbf{Rayleigh waves} (elliptical rolling 
  motion).
\end{itemize}

The propagation velocity of seismic waves is governed by the elastic properties 
and density of the medium. Seismic stations record these ground motions using 
specialized instruments: \textbf{seismometers}, which typically measure ground 
velocity and are sensitive to weak motions, and \textbf{accelerometers}, which 
measure ground acceleration and are essential for recording strong ground 
motions without clipping.

Modern seismic stations continuously record ground motion in three orthogonal 
components: Vertical (Z), North-South (N), and East-West (E). The data are 
digitized and stored in standard formats such as miniSEED (MSEED) or SAC, 
accompanied by metadata (StationXML) containing valid station coordinates and 
instrument response information.